
\begin{titlepage}


\newgeometry{left=2cm, right=1cm, top=1.5cm, bottom=1.5cm}
\centering

\begin{minipage}[c]{0.16\textwidth}
\centering
\includegraphics[width=0.85\textwidth]{\Tapa/figs/unr.jpg}    
\end{minipage}
\begin{minipage}[c]{0.65\textwidth}
\centering{}
{\scshape \small
Departamento de F�sica\\
Escuela de Ciencias Exactas y Naturales\\
Facultad de Ciencias Exactas, Ingenier�a y Agrimensura\\
Universidad Nacional de Rosario\\
}
\end{minipage}
\begin{minipage}[c]{0.16\textwidth}
\centering
\includegraphics[width=1\textwidth]{\Tapa/figs/fceia.jpg}    
\end{minipage}


\vspace*{1.8cm}
\centering

{\large \bf TESINA DE GRADO PARA OPTAR POR\\}
{\large \bf EL T�TULO DE LICENCIADO EN F�SICA\\}

\vspace*{1\baselineskip} % White space at the top of the page


\begin{minipage}[c]{0.9\textwidth}
    \centering
    
    \rule{\textwidth}{1.6pt}\vspace*{-\baselineskip}\vspace*{2pt} % Thick horizontal line
    \rule{\textwidth}{0.4pt}\\[\baselineskip] % Thin horizontal line
    
    \begin{spacing}{2.0}
    {\huge \bf  M�todos Espectrales para la Determinaci�n de la Profundidad del Punto de Curie y el Espesor El�stico de la Corteza Terrestre}\\[0.2\baselineskip] % Title
    \end{spacing}
    
    \rule{\textwidth}{0.4pt}\vspace*{-\baselineskip}\vspace{3.2pt} % Thin horizontal line
    \rule{\textwidth}{1.6pt}\\[\baselineskip] % Thick horizontal line
\end{minipage}

\centering
{\scshape % Texto en mat�sculas con primera letra grande y el resto m�s chicas

\vspace{0.5cm}

Autor\\
{\Large Santiago Rub�n Soler \par} % Editor list
\vspace{0.5cm}

Director\\
{\Large Dr. Mario E. Gim�nez \par} % Editor list

\vspace{0.5cm}

Consejero de Estudios\\
{\Large Dr. Pablo A. Turner \par} % Editor list
}




\vfill


\includegraphics[width=0.14\textwidth]{\Tapa/figs/igsv.png}

\vspace{0.1cm}
{\scshape \small En colaboraci�n con el \\ \small Instituto Geof�sico-Sismol�gico Ing. Volponi \\ \small Universidad Nacional de San Juan}

\vspace{0.5cm}
{\scshape \large 2015}

\restoregeometry
\end{titlepage}


